% paper.tex — LLNCS (ApplePies: Applications in Electronics Pervading Industry,
%             Environment and Society; Springer LNEE)
% Compile: pdflatex paper && bibtex paper && pdflatex paper && pdflatex paper
\documentclass[runningheads]{llncs}
\usepackage[T1]{fontenc}
\usepackage{graphicx}
\usepackage{booktabs}
\usepackage{tabularx}
\usepackage{amsmath}
\usepackage{xspace}
\setlength{\emergencystretch}{3em}
\newcommand{\uJ}{\,$\mu$J}
\begin{document}

\title{A Byte-Level Multilingual News Recommender for Microcontrollers:
       NAS, Micro-NAS and Binarized Micro-NAS across Flash, RAM and Energy}
\titlerunning{Byte-Level Multilingual News Recommender}
\author{}
\authorrunning{}
\institute{}
\maketitle

\begin{abstract}
On-device news recommendation is private (no behavioural data leaves the device)
but faces a cold-start problem and a hard memory wall: classic recommenders carry
a word-embedding table of tens to hundreds of megabytes, beyond a
microcontroller's flash. We present a language-agnostic \emph{content} encoder
that distils a frozen multilingual sentence-transformer into a byte-level char-CNN
whose entire vocabulary is the 256 UTF-8 bytes, removing per-language tables and
covering all 14 xMIND languages with one sub-2\,MB model. We study three search
regimes (NAS, Micro-NAS, binarized Micro-NAS) across FP32/INT8/Binary, reporting
ranking quality (AUC, MRR, nDCG on MIND) against flash, RAM and an energy proxy.
INT8 Micro-NAS reaches \textbf{0.610} AUC at \textbf{204\,KB}, matching an NRMS
word-embedding baseline (0.607) at ${\sim}1/130$ the size (vs.\ 27\,MB); INT8 NAS
reaches \textbf{0.647} at 790\,KB. Naive binary is weak (0.52), but ReActNet-style
training with distilled initialisation recovers it to \textbf{0.572} at 186\,KB.
Distillation adds $+3.3$ AUC; run-to-run std is $\le0.012$; a 64-d Matryoshka
prefix is nearly free; and the exported population prior lifts cold users from
chance (0.500) to 0.551. Multilingual numbers are cross-lingual transfer through
one content encoder (the user model uses English clicks); on-device Pi\,5/STM32H7
measurements are future work. We release the pipeline and artefacts.
\keywords{Edge AI \and Neural Architecture Search \and Binarized Neural Networks
\and Multilingual News Recommendation \and TinyML \and Knowledge Distillation}
\end{abstract}

\section{Introduction}
Privacy-first readers personalise \emph{on-device}, where the user model never
leaves the node. Two obstacles dominate: \textbf{cold start} (a per-user model
from zero recommends poorly) and the \textbf{memory wall} (NRMS~\cite{wu2019nrms},
NAML~\cite{wu2019naml} encode titles with a word-embedding matrix, e.g.\
GloVe-300d~\cite{pennington2014glove} over 30k--100k words $=$ 37--120\,MB, which
cannot fit an MCU, and multiplies per language). We pretrain offline, on
MIND~\cite{wu2020mind} and its multilingual extension xMIND~\cite{iana2024xmind},
a population cold-start prior and a compact content encoder, then compress to the
edge. The encoder is a \textbf{byte-level char-CNN} distilled
\cite{reimers2020multilingual} from a frozen multilingual teacher: its only
symbols are the 256 UTF-8 bytes, so one model handles all 14 xMIND languages with
no per-language table.

\noindent\textbf{Contributions.} (1)~a language-agnostic byte-level content
encoder that removes the embedding-table memory wall and serves 14 languages with
one sub-2\,MB model; (2)~a systematic architecture$\times$precision benchmark
(NAS / Micro-NAS / binarized-$\mu$NAS at FP32/INT8/Binary) against flash, RAM and
energy; (3)~an improved binary encoder (ReActNet + Bi-Real, distilled) lifting
1-bit AUC $0.521\to0.572$ at 186\,KB, plus multi-seed, Matryoshka and cold-start
analyses and released artefacts. To our knowledge this is the first byte-level,
vocabulary-free news recommender studied under a joint architecture-search and
precision sweep for microcontroller-class targets; prior on-device NAS/BNN work is
vision~\cite{lin2020mcunet,nasbnn2024,pighetti2026nasbnn} and prior multilingual
recommenders~\cite{iana2024xmind} are server-scale.

\section{Related Work}
NRMS~\cite{wu2019nrms}, NAML~\cite{wu2019naml}, LSTUR~\cite{an2019lstur},
NPA~\cite{wu2019npa} and DKN~\cite{wang2018dkn} set attention/knowledge title and
user encoders on MIND; the task is impression-level click ranking (AUC/MRR/nDCG).
Published MIND large-test AUC spans 0.65 (DKN) to 0.68 (NRMS)~\cite{wu2020mind};
PLM-based models (PLM-NR~\cite{wu2021plmnr}, Fastformer~\cite{wu2021fastformer},
UNBERT/MINER) reach $0.70$--$0.73$~\cite{wu2023survey} but wrap BERT-scale
encoders (hundreds of MB), unsuitable for MCUs, so we use NRMS, the strongest
lightweight title-only model, as reference. Byte/char models
\cite{clark2022canine,xue2022byt5} and hashing~\cite{joulin2017fasttext} drop the
vocabulary table; multilingual distillation~\cite{reimers2020multilingual} aligns
translations to a shared space. MCU
NAS~\cite{lin2020mcunet,liberis2021unas} and binary NAS~\cite{nasbnn2024} plus
XNOR-Net~\cite{rastegari2016xnor}, Bi-Real~\cite{liu2018bireal} and
ReActNet~\cite{liu2020reactnet} inform the compression; we bring this to on-device
multilingual news recommendation.

\section{Method}
\textbf{Task.} Impression-level click ranking: score candidates given a user's
clicked history. \textbf{Encoder.} A frozen
\texttt{paraphrase-multilingual-MiniLM-L12-v2}~\cite{wang2020minilm} gives a 384-d
anchor per English title; as xMIND titles are parallel translations, every
language is distilled to the \emph{same} anchor. The student is a byte-level
char-CNN: a 256-value byte embedding, depthwise-separable 1-D
conv blocks (cheap, MCU-friendly), masked mean pooling
and a linear head, trained with a Matryoshka~\cite{kusupati2022matryoshka} cosine
loss so a truncated prefix stays usable. \textbf{Recommender.} News vectors feed
an additive-attention user encoder over history; scores are user$\cdot$candidate
dot products, trained NRMS-style with one positive and $K$ negatives (softmax
cross-entropy, AdamW~\cite{loshchilov2019adamw}). \textbf{Search (three arms).}
One space over $\{$channels,\,depth,\,out\_dim$\}$, scored by distillation quality
under footprint feasibility, via evolutionary search: \emph{NAS} (FP32, loose;
accuracy ceiling), \emph{Micro-NAS} (INT8, hard STM32H7 budget), \emph{binarized
Micro-NAS} (Binary, same budget); precision is the second axis.
\textbf{Precision.} STE fake-quantization: per-channel symmetric INT8 and
sign$\times$scale binary; following XNOR-Net~\cite{rastegari2016xnor} the byte
embedding and first/last layers stay FP. The binary arm adds ReActNet activations
(RSign, RPReLU)~\cite{liu2020reactnet} and Bi-Real shortcuts~\cite{liu2018bireal}
with a distilled initialisation. \textbf{Footprint/energy.} Params/MACs via
\texttt{thop}; size $=$ params$\times$bytes/precision; energy proxy $=$
MACs$\times$Horowitz-45\,nm per-op~\cite{horowitz2014} (FP32 MAC 4.6\,pJ, INT8
0.23\,pJ, 1-bit 0.0072\,pJ), a \emph{compute-only} figure that ignores
memory/SRAM energy, used only for relative comparison.

\section{Experimental Setup}
\textbf{Data.} MINDsmall~\cite{wu2020mind} (50k users; train 156{,}965 / dev
73{,}152 impressions; 51{,}282 news) and xMINDsmall~\cite{iana2024xmind} (14
languages, NLLB-3.3B translations~\cite{nllb2022}). xMIND supplies translated
titles keyed by news id; behaviours are reused. \emph{Every recommender and NRMS
trains on English MINDsmall clicks only} (like-for-like); xMIND is used only for
cross-lingual evaluation and for the distilled-init distillation, so multilingual
results are cross-lingual transfer, not native-language personalisation.
MINDsmall (not MINDlarge) keeps the full study (3 arms, 3 precisions, NAS,
15-language eval) to 8.2\,h on one laptop; small has no public test, so we use
dev and cite NRMS MINDlarge-test 0.6776~\cite{wu2020mind} as external reference.
\textbf{NAS.} Evolutionary, population 30, 15 generations; candidates scored by a
2-epoch distillation on 20k titles, budget-infeasible ones rejected; Micro-NAS
thresholds (256\,KB / 128\,KB / 2\,M MACs) sit ${\sim}8\times$ below the H7
ceiling so they bind. \textbf{Hardware.} Training/search on one laptop: Intel Core
Ultra~9 275HX (24c), RTX~5070 Laptop GPU (Blackwell, 8\,GB), 32\,GB, PyTorch~2.12
(CUDA~13). Edge targets: Raspberry~Pi~5 (Cortex-A76, onnxruntime) and STM32H743
(Cortex-M7 @480\,MHz, 2\,MB/1\,MB, X-CUBE-AI); their on-device latency/energy are
\textbf{future work} (boards pending), so only laptop latency and the analytical
proxy are reported. \textbf{Reproducibility.} Fixed seed; SHA256-pinned data;
code, notebook, artefacts and a dependency lock are released.\footnote{Repository
URL omitted for double-blind review.}

\section{Results}
\begin{table}[t]
\caption{Architecture $\times$ precision (MINDsmall dev). \textbf{Bold}: best per
group; $\dagger$: native precision. \emph{Arm} $=$ search procedure (a distinct
architecture); \emph{precision} $=$ deployment quantization; off-diagonal cells
are ablations. Energy is the Horowitz proxy.}
\label{tab:matrix}\centering\footnotesize\setlength{\tabcolsep}{4pt}
\begin{tabular}{llccccr}
\toprule
Arm & Prec. & AUC & MRR & nDCG@10 & Size\,(KB) & Energy\,($\mu$J)\\
\midrule
NAS & FP32$^\dagger$ & 0.633 & 0.326 & 0.375 & 1566.8 & 168.7\\
NAS & INT8 & \textbf{0.647} & \textbf{0.338} & \textbf{0.387} & 789.8 & 4.65\\
NAS & Binary & 0.588 & 0.297 & 0.342 & 563.1 & 4.25\\
\midrule
Micro-NAS & FP32 & 0.605 & 0.307 & 0.352 & 266.8 & 15.90\\
Micro-NAS & INT8$^\dagger$ & \textbf{0.610} & \textbf{0.314} & \textbf{0.360} & 203.9 & 2.30\\
Micro-NAS & Binary & 0.521 & 0.251 & 0.294 & 185.6 & 2.25\\
\midrule
Bin.\,$\mu$NAS & FP32 & 0.593 & 0.290 & 0.336 & 311.4 & 15.43\\
Bin.\,$\mu$NAS & INT8 & 0.601 & 0.295 & 0.344 & 255.7 & 3.28\\
Bin.\,$\mu$NAS & Binary$^\dagger$ & \textbf{0.546} & \textbf{0.273} & \textbf{0.317} & 239.5 & 3.23\\
\bottomrule
\end{tabular}
\end{table}

Table~\ref{tab:matrix} (arms NAS 256-4-384, Micro-NAS 64-5-384, Bin.\,$\mu$NAS
96-2-384) shows INT8 is effectively lossless vs.\ FP32 while cutting energy
7--36$\times$ and size $\approx1.4$--$2\times$; binarization costs 5--9 AUC points
but, since the embedding and first/last layers stay FP, its size gain over INT8 is
modest here. Single-inference laptop latency (batch~1) is 0.34/0.19/0.10\,ms on
CPU (onnxruntime) for NAS/Micro-NAS/Bin.\,$\mu$NAS and $<0.7$\,ms on GPU --- all
far under a 50\,ms budget; CPU latency tracks MACs.

\noindent\textbf{Baseline.} Our reproduced NRMS (256-d, 25{,}355-word vocabulary)
has 7.08\,M parameters, 91.6\% in the embedding table ($\sim$27\,MB FP32; INT8 is
6.8\,MB with $\Delta$AUC $<0.001$), reaching dev AUC 0.607. The byte student
\emph{matches} it at 204\,KB (33$\times$ smaller than even NRMS-INT8) and
\emph{exceeds} it at 790\,KB (INT8 NAS, 0.647), with no vocabulary table.

\noindent\textbf{Multilingual (cross-lingual transfer).} Per-language dev AUC of
Micro-NAS/INT8, one byte model, no per-language table: en~.604, hat~.556,
jpn~.551, ind~.537, grn~.533, ron~.532, tur~.530, vie~.523, swh~.521, zho~.521,
som~.514, tha~.504, tam~.504, fin~.495, kat~.491. Transfer is strongest for
Latin-script / higher-resource targets and weakest for distinct-script,
low-resource languages (kat, fin), tracking NLLB quality. These measure the
content encoder's cross-lingual reach, not native-language personalisation.

\noindent\textbf{Distillation ablation.} At Micro-NAS (64-5-384, 10 epochs),
initialising the encoder from the distilled student rather than training from
scratch raises AUC $0.600\to0.633$ (MRR $0.293\to0.345$), so the gain comes from
distillation, not capacity; we deploy the distilled-init encoder.

\noindent\textbf{Binary recovery.} Naive 1-bit is near chance (Micro-NAS/Binary
0.521); ReActNet + Bi-Real with distilled init lifts it to \textbf{0.572} at
186\,KB, within 0.035 AUC of the FP32 baseline at two orders of magnitude less
memory --- a usable ultra-low-footprint option, though INT8 leads on accuracy.

\noindent\textbf{Robustness, Matryoshka, cold start.} Over 3 seeds Micro-NAS/INT8
is $0.624\pm0.004$ and std stays $\le0.012$, so single-seed numbers are stable.
Truncating the student embedding $384\to64$ dims barely changes cosine to the
teacher ($0.454\to0.459$), so an on-device 64-d embedding is nearly free. With
history masked, the history-only model is at chance (0.500); ranking cold users by
the exported topic prior lifts them to 0.551, confirming its value.

\noindent\textbf{Trade-off.} On the accuracy--footprint front, Micro-NAS/INT8
(0.610, 204\,KB, 2.3\uJ) is the on-device sweet spot and NAS/INT8 the accuracy
leader (0.647, 790\,KB); binary is smallest but footprint-motivated only.

\section{Limitations}
\textbf{On-device measurements are pending} (Pi\,5/STM32H7 boards unavailable): we
report laptop latency, computed size/RAM, and the Horowitz energy proxy, which
ignores SRAM/memory traffic. \textbf{Binary deployment is analytical}: INT8
deploys via ONNX$\to$X-CUBE-AI, but the 1-bit path needs Larq/QKeras (TensorFlow;
Larq is archived and its compute engine has no Windows build), and Cortex-M7 lacks
POPCOUNT, so software bit-counting can erase the MAC saving --- we make no binary
\emph{speed} claim. \textbf{Byte sequences} are ${\sim}5\times$ longer than words,
trading flash for on-MCU SRAM bandwidth (to be quantified on-device).
\textbf{Cross-lingual transfer}, not native behaviour; \textbf{MINDsmall} scale;
\textbf{MT quality} varies for low-resource languages.

\section{Ethics and Conclusion}
MIND (Microsoft Research License) and xMIND (CC-BY-NC-SA-4.0) are used for
non-commercial research; no personal data is collected; translations are
machine-generated. A byte-level language-agnostic encoder removes the
embedding-table memory wall and serves 14 languages with one sub-2\,MB model.
INT8 Micro-NAS is the practical sweet spot (204\,KB, 2.3\uJ, AUC 0.610); binary,
recovered to 0.572 at 186\,KB, is the ultra-low-footprint option, with real 1-bit
ARM execution (Larq/X-CUBE-AI on Linux) and on-device measurement as the main
future work. Pipeline and artefacts (ONNX encoder, cold-start prior) are released.

\bibliographystyle{splncs04}
\bibliography{references}
\end{document}
